\section{Mathematical Model}
\label{sec:model}

The 3D-SCPTOP formulation integrates stack-based truck-loading
constraints with procurement transport and inventory decisions. It
departs from a standard 3D container-loading formulation in two
related ways. First, the 3D placement of individual RTIs inside the
truck is replaced by a 2D placement of variable-composition
\emph{stacks} on the truck floor, with per-stack height and weight
caps. Second, the per-period decision of which product fills each slot
is retained, but the crush-limit cap is evaluated slot-by-slot using
the period's actual product mix. This yields a tighter per-stack
weight bound than a worst-case-per-RTI-type abstraction, while
eliminating vertical-stability constraints altogether, since stacks
are self-supporting by construction.

The model is written over \emph{trucks}. A dispatched truck designs its
own layout: the stack types, orientations, positions and layer
compositions of its slots are variables of that truck and of that
period, and the load it carries is assigned to those slots in the same
period. There is no catalogue of loading configurations, and no
variable that assigns a truck to an entry of one. This is a modelling
decision and not a presentational one, and it is the decision Section
\ref{sec:math} rests on: every row below carries either the triple
$(g,k,t)$ or none of it, so the model is block-angular over
$G\times K\times T$ and the pricing problem of the Dantzig--Wolfe
reformulation is obtained by deleting that quantifier and nothing else.
Section \ref{sec:decomp} performs exactly that deletion, row by row.

\paragraph{Notational conventions.}
Throughout this section, $g\in G=\{s,p\}$ indexes nodes (supplier and
plant), $n\in N$ products, $r\in R$ RTI types, $h\in H$ stack types,
$i,j\in I$ stack-slot positions on a truck floor, $k\in K$ trucks and
$t\in T$ planning periods. Each product travels in exactly one RTI
type, written $\rho(n)$, and $M(r)=\{n\in N:\rho(n)=r\}$ is the family
of products that RTI type $r$ carries. Variables whose direction is
determined by definition omit the $g$ subscript: $V_{kint}$ denotes
full RTIs in slot $i$ of a plant-bound truck, and $V^{\mathrm{e}}_{krt}$
folded RTIs on a supplier-bound truck. Table
\ref{tab:sb-Nomenclature} collects the complete nomenclature.

% -----------------------------------------------------------
\begin{longtable}[hbt!]{|p{0.15\textwidth}|p{0.79\textwidth}|}
    \caption{Nomenclature of 3D-SCPTOP (truck-indexed stack-based formulation)}
    \label{tab:sb-Nomenclature}\\
    \hline
    \multicolumn{2}{|l|}{\textbf{Indices and sets}}\\\hline
    $G$            & Supplier and plant nodes, $G=\{s,p\}$; $g=s$ is the folded
                     direction (plant to supplier), $g=p$ the full one.\\
    $N$            & Products, $n=1,\dots,|N|$.\\
    $R$            & RTI types, $r=1,\dots,|R|$.\\
    $H$            & Stack types, $h=1,\dots,|H|$.\\
    $I$            & Stack-slot indices on a truck floor, $i=1,\dots,|I|$ (see
                     \eqref{sb:eq:Ibound}).\\
    $K$            & Trucks available per direction and period, $k=1,\dots,|K|$.\\
    $T$            & Planning periods, $t=1,\dots,|T|$.\\
    $\rho(n)$      & The RTI type in which product $n$ travels.\\
    $M(r)$         & Products carried by RTI type $r$, $M(r)=\{n:\rho(n)=r\}\subseteq N$.\\
    $R(h)$         & RTI types compatible with stack type $h$, $R(h)\subseteq R$.\\
    \hline
    \multicolumn{2}{|l|}{\textbf{Parameters}}\\\hline
    $cK$           & Cost of dispatching one truck.\\
    $lt$           & Lead transport time between plant and supplier (and vice versa), $lt\ge1$.\\
    $i_{gn}$       & Initial inventory of $n\in N$ at $g\in G$ (full RTIs).\\
    $ie_{gr}$      & Initial inventory of folded RTIs $r\in R$ at $g\in G$.\\
    $ic_{gn}$      & Inventory holding cost of $n\in N$ at $g\in G$.\\
    $d_{gnt}$      & Production plan for $n\in N$ at $g\in G$ in $t\in T$: filled at $s$, consumed at $p$.\\
    $a_{nt}$       & Pipeline arrivals of full RTIs of product $n$ at $p$ in $t\in\{1,\dots,lt\}$.\\
    $ae_{rt}$      & Pipeline arrivals of folded RTIs of type $r$ at $s$ in $t\in\{1,\dots,lt\}$.\\
    $ss_{gnt}$     & Coverage stock of $n$ at $g$ in $t$ (full RTIs).\\
    $sse_{gtr}$    & Coverage stock of folded RTIs $r$ at $g$ in $t$.\\
    $w_{n}$        & Weight of a full RTI carrying product $n$; $w_{n}>0$.\\
    $w^{\mathrm{f}}_{r}$ & Weight of a folded RTI of type $r$; $w^{\mathrm{f}}_{r}>0$.\\
    $pK^{w}$       & Truck payload.\\
    $W^{\max}$     & Per-stack crush-limit weight cap: the \emph{total} weight of one stack, summed over all its layers, not the weight borne by its bottom RTI alone.\\
    $pK,qK,rK$     & Truck floor length, floor width, and inner height.\\
    $\Gamma_K$     & Truck volume, $\Gamma_K=pK\cdot qK\cdot rK$.\\
    $\alpha$       & Minimum truck utilisation, required on whichever of volume or payload binds ($0<\alpha\le 1$).\\
    $p_{h},q_{h}$  & Footprint length and width of stack type $h$.\\
    $r_{gr}$       & Height of one RTI of type $r$ at node $g$ (full at $p$, folded at $s$); $r_{gr}>0$.\\
    $a_{r}$        & Footprint area of one RTI of type $r$; $a_{r}>0$.\\
    $\gamma_{gr}$  & Volume of one RTI of type $r$ in direction $g$, $\gamma_{gr}=a_{r}\,r_{gr}$.\\
    $n_{hr}$       & RTIs of type $r\in R(h)$ per layer of stack type $h$; one layer tiles the footprint, $n_{hr}a_{r}=p_{h}q_{h}$.\\
    $\omega_{gr}$  & Weight per RTI used by the layer-capacity form of the
                     utilisation floor: $\omega_{sr}=w^{\mathrm{f}}_{r}$ and
                     $\omega_{pr}=\max_{n\in M(r)}w_{n}$.\\
    $\lambda^{K},\lambda^{I}$ & Weights of the transport and inventory objectives; both $1$ throughout this paper.\\
    $M^{x},M^{y}$  & Big-$M$ constants of the non-overlap rows, $M^{x}=pK$ and $M^{y}=qK$.\\
    $\bar u^{V}_{n}$ & Most full RTIs of $n$ one stack can carry, $\bar u^{V}_{n}=\lfloor W^{\max}/w_{n}\rfloor$.\\
    $\bar u^{\mathrm{e}}_{r}$ & Most folded RTIs of $r$ one truck can carry, $\bar u^{\mathrm{e}}_{r}=\lfloor pK^{w}/w^{\mathrm{f}}_{r}\rfloor$.\\
    \hline
    \multicolumn{2}{|l|}{\textbf{Decision variables -- the truck and its layout}}\\\hline
    $Z_{gkt}$        & $\in\{0,1\}$. Equals $1$ iff truck $k$ of direction $g$ in period $t$ is configured, i.e.\ carries a layout.\\
    $Y_{gkt}$        & $\in\{0,1\}$. Equals $1$ iff truck $k$ is dispatched in direction $g$ in period $t$.\\
    $\theta_{gkt}$   & $\in\{0,1\}$. Selects the resource on which the utilisation floor of that truck is required: volume if $1$, payload if $0$.\\
    $X_{gihkt}$      & $\in\{0,1\}$. Equals $1$ iff slot $i$ of that truck holds a stack of type $h$.\\
    $L_{gihktr}$     & $\in\mathbb{Z}_{\ge 0}$. Layers of RTI type $r\in R(h)$ in slot $i$ of that truck.\\
    $LX_{gihkt}$     & $\in\{0,1\}$. Equals $1$ iff stack $h$'s length lies along the $x$-axis in slot $i$.\\
    $LY_{gihkt}$     & $\in\{0,1\}$. Equals $1$ iff stack $h$'s length lies along the $y$-axis in slot $i$.\\
    $PX_{gikt},PY_{gikt}$ & $\ge 0$. Bottom-left coordinates of slot $i$ on the truck floor.\\
    $A_{gijkt}$      & $\in\{0,1\}$. Equals $1$ iff slot $i$ is to the left of slot $j$.\\
    $B_{gijkt}$      & $\in\{0,1\}$. Equals $1$ iff slot $i$ is to the right of slot $j$.\\
    $C_{gijkt}$      & $\in\{0,1\}$. Equals $1$ iff slot $i$ is behind slot $j$.\\
    $D_{gijkt}$      & $\in\{0,1\}$. Equals $1$ iff slot $i$ is in front of slot $j$.\\
    \hline
    \multicolumn{2}{|l|}{\textbf{Decision variables -- load and inventory}}\\\hline
    $V_{kint}$       & $\in\mathbb{Z}_{\ge 0}$. Full RTIs of product $n$ in slot $i$ of plant-bound truck $k$ in period $t$.\\
    $V^{\mathrm{e}}_{krt}$ & $\in\mathbb{Z}_{\ge 0}$. Folded RTIs of type $r$ on supplier-bound truck $k$ in period $t$.\\
    $I_{gnt}$        & $\in\mathbb{Z}_{\ge 0}$. End-of-period inventory of $n$ at $g$ in $t$ (full RTIs).\\
    $Ie_{gtr}$       & $\in\mathbb{Z}_{\ge 0}$. End-of-period inventory of folded RTIs $r$ at $g$ in $t$.\\
    \hline
    \multicolumn{2}{|l|}{\textbf{Derived quantities} (abbreviations, not constraints; see \eqref{sb:eq:agg-V}--\eqref{sb:eq:agg-L})}\\\hline
    $\overline{V}_{nt}$ & Full RTIs of $n$ delivered to the plant in $t$.\\
    $\overline{V^{\mathrm{e}}}_{rt}$ & Folded RTIs of $r$ returned to the supplier in $t$.\\
    $\overline{V}_{knt}$ & Full RTIs of $n$ on plant-bound truck $k$ in $t$.\\
    $\widehat{X}_{gikt}$ & Equals $1$ iff slot $i$ of that truck is used.\\
    $\overline{L}_{ghktr}$ & Layers of RTI type $r$ under stack type $h$ on that truck, over all slots.\\
    \hline
\end{longtable}

% --------------------------
% Derived quantities
% --------------------------
\subsection*{Derived quantities}

Five sums recur, and naming them is what keeps the model readable and
what makes its structure visible. They are \emph{abbreviations}, not
constraints: each is a linear expression in the decision variables
above, substituted wherever it appears.
\begin{align}
\overline{V}_{nt}&=\sum_{k\in K}\sum_{i\in I} V_{kint}
    &&\forall n\in N,\; t\in T
    \label{sb:eq:agg-V}\\
\overline{V^{\mathrm{e}}}_{rt}&=\sum_{k\in K} V^{\mathrm{e}}_{krt}
    &&\forall r\in R,\; t\in T
    \label{sb:eq:agg-Ve}\\
\overline{V}_{knt}&=\sum_{i\in I} V_{kint}
    &&\forall k\in K,\; n\in N,\; t\in T
    \label{sb:eq:agg-Vk}\\
\widehat{X}_{gikt}&=\sum_{h\in H} X_{gihkt}
    &&\forall g\in G,\; i\in I,\; k\in K,\; t\in T
    \label{sb:eq:agg-X}\\
\overline{L}_{ghktr}&=\sum_{i\in I} L_{gihktr}
    &&\forall g\in G,\; h\in H,\; k\in K,\; t\in T,\; r\in R(h)
    \label{sb:eq:agg-L}
\end{align}

The first two are the hinge of the whole paper. The inventory block
reads $\overline{V}_{nt}$ and $\overline{V^{\mathrm{e}}}_{rt}$ and
nothing else about the trucks: it does not see which truck carried
what, nor how the load was placed. Every other row of the model carries
a single $(g,k,t)$. The entire coupling between the transport block and
the loading block is therefore the two sums
\eqref{sb:eq:agg-V}--\eqref{sb:eq:agg-Ve} together with the fact that
there are at most $|K|$ trucks per direction and period --- which is
why the master of Section \ref{sec:mp} needs exactly two definition
rows and one count row, and why
\eqref{sb:eq:agg-V}--\eqref{sb:eq:agg-Ve} reappear there as the only
place a column's image enters.

% --------------------------
% Objective
% --------------------------
\subsection*{Objectives}

The model has two objectives, transport cost and inventory holding
cost:
\begin{align}
\min\; z^{K} &= \sum_{g\in G}\sum_{k\in K}\sum_{t\in T} cK\cdot Y_{gkt}
    \label{sb:eq:obj-zK}\\
\min\; z^{I} &= \sum_{g\in G}\sum_{n\in N}\sum_{t\in T} ic_{gn}\cdot I_{gnt}
    \label{sb:eq:obj-zI}
\end{align}
and they are solved scalarised,
\begin{equation}
\min\;\; \lambda^{K} z^{K}+\lambda^{I} z^{I},
\label{sb:eq:obj}
\end{equation}
with $\lambda^{K}=\lambda^{I}=1$ in every run reported in Section
\ref{sec:results}, so that the reported cost is transport plus holding
cost in the same money. Transport cost counts \emph{dispatches}: a
truck that is configured but not dispatched costs nothing, and no term
of \eqref{sb:eq:obj-zK}--\eqref{sb:eq:obj-zI} charges for a layout.

% --------------------------
% Inventory balance
% --------------------------
\subsection*{Inventory and flow balance}

Full RTIs arrive at the plant $lt$ periods after they are dispatched;
inventory itself accumulates one period at a time, and only the arrival
is delayed. For the plant,
\begin{align}
I_{pnt}  &=I_{p,n,t-1}-d_{pnt}+\overline{V}_{n,t-lt}
    &&\forall n\in N,\; t\ge lt+1
    \label{sb:eq:inv-plant-main}\\
I_{pnt}  &=I_{p,n,t-1}-d_{pnt}+a_{nt}
    &&\forall n\in N,\; t\in\{2,\dots,lt\}
    \label{sb:eq:inv-plant-transient}\\
I_{pnt}  &=i_{pn}-d_{pnt}+a_{nt}
    &&\forall n\in N,\; t=1
    \label{sb:eq:inv-plant-init}\\
Ie_{ptr} &=Ie_{p,t-1,r}+\sum_{n\in M(r)} d_{pnt}-\overline{V^{\mathrm{e}}}_{rt}
    &&\forall r\in R,\; t\ge 2
    \label{sb:eq:inve-plant}\\
Ie_{ptr} &=ie_{pr}+\sum_{n\in M(r)} d_{pnt}-\overline{V^{\mathrm{e}}}_{rt}
    &&\forall r\in R,\; t=1
    \label{sb:eq:inve-plant-init}
\end{align}
Consuming a full RTI at the plant produces a folded one in the same
period, which is why the plant's folded balance is driven by $d_{pnt}$;
the folded RTIs leave on the same day they are shipped, so
$\overline{V^{\mathrm{e}}}_{rt}$ enters \eqref{sb:eq:inve-plant}
undelayed and arrives at the supplier $lt$ periods later. For the
supplier,
\begin{align}
I_{snt}  &=I_{s,n,t-1}+d_{snt}-\overline{V}_{nt}
    &&\forall n\in N,\; t\ge 2
    \label{sb:eq:inv-sup-main}\\
I_{snt}  &=i_{sn}+d_{snt}-\overline{V}_{nt}
    &&\forall n\in N,\; t=1
    \label{sb:eq:inv-sup-init}\\
Ie_{str} &=Ie_{s,t-1,r}+\overline{V^{\mathrm{e}}}_{r,t-lt}-\sum_{n\in M(r)} d_{snt}
    &&\forall r\in R,\; t\ge lt+1
    \label{sb:eq:inve-sup-main}\\
Ie_{str} &=Ie_{s,t-1,r}+ae_{rt}-\sum_{n\in M(r)} d_{snt}
    &&\forall r\in R,\; t\in\{2,\dots,lt\}
    \label{sb:eq:inve-sup-transient}\\
Ie_{str} &=ie_{sr}+ae_{rt}-\sum_{n\in M(r)} d_{snt}
    &&\forall r\in R,\; t=1
    \label{sb:eq:inve-sup-init}
\end{align}
Coverage stock is enforced at both nodes and on both kinds of stock:
\begin{align}
I_{gnt}&\ge ss_{gnt}
    &&\forall g\in G,\; n\in N,\; t\in T
    \label{sb:eq:safety-full}\\
Ie_{gtr}&\ge sse_{gtr}
    &&\forall g\in G,\; r\in R,\; t\in T
    \label{sb:eq:safety-empty}
\end{align}

These are the only rows of the model that couple two trucks, two
periods or two directions, and they do it only through
\eqref{sb:eq:agg-V}--\eqref{sb:eq:agg-Ve}.

% --------------------------
% Stack packing block
% --------------------------
\subsection*{The layout of a truck}

Everything from here to \eqref{sb:eq:stack-w-sup} is written for one
truck: one direction $g$, one truck index $k$, one period $t$. At most
one stack type occupies a slot, and only a configured truck has slots:
\begin{equation}
\widehat{X}_{gikt}\;\le\;Z_{gkt}
\qquad \forall g\in G,\; i\in I,\; k\in K,\; t\in T.
\label{sb:eq:slot-active}
\end{equation}

The stack occupying slot $i$ is oriented either with its length along
the $x$-axis or along the $y$-axis:
\begin{equation}
LX_{gihkt}+LY_{gihkt}\;=\;X_{gihkt}
\qquad \forall g\in G,\; h\in H,\; i\in I,\; k\in K,\; t\in T.
\label{sb:eq:orient}
\end{equation}

The total stack height cannot exceed the truck height. The following
constraint also forces $L_{gihktr}=0$ whenever $X_{gihkt}=0$, since
$r_{gr}>0$:
\begin{equation}
\sum_{r\in R(h)} r_{gr}\,L_{gihktr}\;\le\;rK\cdot X_{gihkt}
\qquad \forall g\in G,\; h\in H,\; i\in I,\; k\in K,\; t\in T.
\label{sb:eq:height}
\end{equation}

Slots must not overlap. For each pair $i<j$ of slots of the same truck,
with $\forall g\in G,\;k\in K,\;t\in T$ understood throughout:
\begin{align}
PX_{gikt}+\sum_{h\in H}\bigl(p_{h}LX_{gihkt}+q_{h}LY_{gihkt}\bigr)
    &\le PX_{gjkt}+M^{x}(1-A_{gijkt})
    \label{sb:eq:nonov-x1}\\
PX_{gjkt}+\sum_{h\in H}\bigl(p_{h}LX_{gjhkt}+q_{h}LY_{gjhkt}\bigr)
    &\le PX_{gikt}+M^{x}(1-B_{gijkt})
    \label{sb:eq:nonov-x2}\\
PY_{gikt}+\sum_{h\in H}\bigl(q_{h}LX_{gihkt}+p_{h}LY_{gihkt}\bigr)
    &\le PY_{gjkt}+M^{y}(1-C_{gijkt})
    \label{sb:eq:nonov-y1}\\
PY_{gjkt}+\sum_{h\in H}\bigl(q_{h}LX_{gjhkt}+p_{h}LY_{gjhkt}\bigr)
    &\le PY_{gikt}+M^{y}(1-D_{gijkt})
    \label{sb:eq:nonov-y2}
\end{align}
with the cover constraint ensuring that whenever both slots are used
at least one relative-position variable is active:
\begin{equation}
A_{gijkt}+B_{gijkt}+C_{gijkt}+D_{gijkt}
\;\ge\;\widehat{X}_{gikt}+\widehat{X}_{gjkt}-1
\qquad \forall g,\,k,\,t,\,i<j.
\label{sb:eq:cover}
\end{equation}

Containment in the truck floor:
\begin{align}
PX_{gikt}+\sum_{h\in H}\bigl(p_{h}LX_{gihkt}+q_{h}LY_{gihkt}\bigr)
    &\le pK
    &&\forall g,i,k,t,
    \label{sb:eq:cont-x}\\
PY_{gikt}+\sum_{h\in H}\bigl(q_{h}LX_{gihkt}+p_{h}LY_{gihkt}\bigr)
    &\le qK
    &&\forall g,i,k,t.
    \label{sb:eq:cont-y}
\end{align}
These two rows are unconditional, and the non-overlap rows take the
tightened constants $M^{x}=pK$ and $M^{y}=qK$, both because the
geometry is now the truck's own. An unused slot has
$LX_{gihkt}=LY_{gihkt}=0$ for every $h$, so
\eqref{sb:eq:cont-x}--\eqref{sb:eq:cont-y} reduce to $PX_{gikt}\le pK$
and $PY_{gikt}\le qK$ there and cut nothing; and given those bounds an
inactive disjunct of \eqref{sb:eq:nonov-x1}--\eqref{sb:eq:nonov-y2} is
never binding, since the left-hand side is at most $pK$ (respectively
$qK$) and the other slot's origin is non-negative. Enlarging the two
constants only weakens the linear relaxation.

% --------------------------
% Runtime full-truckload
% --------------------------
\subsection*{Full-truckload on every dispatched truck}

Full-truckload transport requires a dispatched truck to be full, but a
truck can be full in two different ways: it may run out of space or it
may run out of payload. Imposing the requirement on volume alone is not
a conservative simplification, it is an infeasible one --- filling the
fraction $\alpha=0.55$ of a $97.92$~m$^{3}$ trailer
($13\,600\times2\,400\times3\,000$~mm) with freight at $579$~kg/m$^{3}$
weighs $31.2$~t against a $24$~t cap, so no truck is dispatchable at all
and the instance is empty. That density is not a worst case: it is
Container~4 of the industrial catalogue of Section \ref{sec:case}
($500$~kg of payload in $1\,200\times800\times900$~mm, tare excluded),
and seven of the catalogue's fifteen types are denser still. Summing the two ratios
instead is the opposite error: it accepts a truck that is $35\%$ full
by volume and $20\%$ by weight, which is full by neither. We therefore
require a dispatched truck to reach a fraction $\alpha$ of
\emph{either} resource, using a binary $\theta_{gkt}$ to select which
one binds. The requirement is stated on the load the truck actually
carries, with each commodity's own weight:
\begin{align}
\sum_{n\in N}\gamma_{p\rho(n)}\,\overline{V}_{knt}
&\;\ge\;\alpha\,\Gamma_{K}\bigl(Y_{pkt}-1+\theta_{pkt}\bigr)
    &&\forall k\in K,\; t\in T
    \label{sb:eq:util-p}\\
\sum_{n\in N} w_{n}\,\overline{V}_{knt}
&\;\ge\;\alpha\,pK^{w}\bigl(Y_{pkt}-\theta_{pkt}\bigr)
    &&\forall k\in K,\; t\in T
    \label{sb:eq:utilw-p}\\
\sum_{r\in R}\gamma_{sr}\,V^{\mathrm{e}}_{krt}
&\;\ge\;\alpha\,\Gamma_{K}\bigl(Y_{skt}-1+\theta_{skt}\bigr)
    &&\forall k\in K,\; t\in T
    \label{sb:eq:util-s}\\
\sum_{r\in R} w^{\mathrm{f}}_{r}\,V^{\mathrm{e}}_{krt}
&\;\ge\;\alpha\,pK^{w}\bigl(Y_{skt}-\theta_{skt}\bigr)
    &&\forall k\in K,\; t\in T
    \label{sb:eq:utilw-s}
\end{align}
At $Y_{gkt}=0$ both rows of a direction are vacuous; at $Y_{gkt}=1$
exactly one is active and the other has a right-hand side of zero and is
vacuous too, so each pair is the disjunction ``full by volume
\emph{or} full by payload''.

The same disjunction can be written on the layer \emph{capacity} the
truck's stacks provide rather than on the load they receive, and in
that form it is a valid tightening. It is gated on the dispatch
$Y_{gkt}$ and not on the configuration $Z_{gkt}$, because it is a
dispatched truck that must be full: a configured, undispatched truck
carries nothing by
\eqref{fix:slot-link}--\eqref{fix:empty-link} and demanding that its
unused layers be full would be a condition on no freight. The
distinction disappears in Section \ref{sec:math}, where the two
binaries are identified.
\begin{align}
\sum_{h\in H}\sum_{r\in R(h)}
    \gamma_{gr}\,n_{hr}\,\overline{L}_{ghktr}
&\;\ge\;\alpha\,\Gamma_{K}\bigl(Y_{gkt}-1+\theta_{gkt}\bigr)
    &&\forall g\in G,\; k\in K,\; t\in T
    \label{sb:eq:util-cap}\\
\sum_{h\in H}\sum_{r\in R(h)}
    \omega_{gr}\,n_{hr}\,\overline{L}_{ghktr}
&\;\ge\;\alpha\,pK^{w}\bigl(Y_{gkt}-\theta_{gkt}\bigr)
    &&\forall g\in G,\; k\in K,\; t\in T
    \label{sb:eq:utilw-cap}
\end{align}
Which of the two forms is the \emph{requirement} and which is the
consequence matters, and the paper states it once here. The requirement
is \eqref{sb:eq:util-p}--\eqref{sb:eq:utilw-s}:
full-truckload is a property of the freight that moves, and a truck
whose stacks \emph{could} have been filled to the floor but were not is
not full. \eqref{sb:eq:util-cap}--\eqref{sb:eq:utilw-cap} are implied
by it --- a slot never receives more RTIs of type $r$ than its layers
hold, by \eqref{sb:eq:slot-cap}--\eqref{sb:eq:empty-cap} below, and
$\gamma_{gr}$ depends only on $r$ while $\omega_{gr}\ge w_{n}$ for every
$n\in M(r)$ --- so they cut nothing and are kept only because they let a
solver establish a first feasible layout without first choosing a
product assignment. Taken \emph{alone}, in place of the load rows, they
would be strictly weaker and would admit trucks that are not full.

The coefficient $\omega_{pr}$ in \eqref{sb:eq:utilw-cap} is the
\emph{heaviest} product the family can carry, and both of the reasons
point the same way. The row is a $\ge$ requirement, so the smaller
coefficient is the stronger row: with $\min_{n\in M(r)}w_{n}$ the row
would assert that \emph{every} product assignment of this layout fills
the truck, which is not a necessary condition of ``this truck is full''
and cuts off loads the problem allows, so that the formulation reports
an optimum above the true one --- in the very comparison where the
decomposition is claimed to beat it. With the maximum the row is
implied by the load rows, as the previous paragraph needs it to be;
with the minimum it is not implied at all. The same principle governs
the crush cap \eqref{sb:eq:stack-w-plant}, which is why that row is
written on the mix actually loaded. A per-load row must use the load's
own product weights; only a row that does not yet know the product may
use a family bound, and then only the bound that keeps it necessary
rather than sufficient.

% --------------------------
% Dispatch, loading and weight
% --------------------------
\subsection*{Dispatch, loading and weight}

A dispatched truck must be configured:
\begin{equation}
Y_{gkt}\;\le\;Z_{gkt}
\qquad \forall g\in G,\; k\in K,\; t\in T.
\label{sb:eq:lc-active}
\end{equation}
The converse does not hold, and the two binaries are not the same
variable in this section: $Z_{gkt}=1$ with $Y_{gkt}=0$ is a truck that
carries a layout and no freight. Such a point costs nothing by
\eqref{sb:eq:obj-zK} and delivers nothing by
\eqref{fix:slot-link}--\eqref{fix:empty-link}; Section \ref{sec:pp}
returns to it, because in the reformulated block it is exactly the
empty column and the two binaries can then be identified.

For plant-bound trucks, full RTIs of type $r$ loaded in slot $i$ cannot
exceed the layer capacity of the stack placed in that slot; products
sharing the same RTI type partition that slot capacity:
\begin{equation}
\sum_{n\in M(r)} V_{kint}
\;\le\;\sum_{h\,:\,r\in R(h)} n_{hr}\,L_{p,i,h,k,t,r}
\qquad \forall k\in K,\; i\in I,\; r\in R,\; t\in T.
\label{sb:eq:slot-cap}
\end{equation}
For supplier-bound trucks, folded RTIs of type $r$ are capped by the
aggregate layer capacity of the truck (no slot-level tracking is needed
because folded-RTI weight is deterministic in $r$):
\begin{equation}
V^{\mathrm{e}}_{krt}
\;\le\;\sum_{h\,:\,r\in R(h)} n_{hr}\,\overline{L}_{s,h,k,t,r}
\qquad \forall k\in K,\; r\in R,\; t\in T.
\label{sb:eq:empty-cap}
\end{equation}
Both rows are unconditional: the layers they refer to are the truck's
own, and a truck with no stacks has zero capacity on the right-hand
side.

Capacity is not the same as dispatch, and the two rows above would
leave the load of a configured but undispatched truck free. The loads
are therefore also gated on the dispatch, at their own analytic bounds
rather than at a generic big-$M$:
\begin{align}
V_{kint}&\le \bar u^{V}_{n}\,Y_{pkt}
    &&\forall k\in K,\; i\in I,\; n\in N,\; t\in T
    \label{fix:slot-link}\\
V^{\mathrm{e}}_{krt}&\le \bar u^{\mathrm{e}}_{r}\,Y_{skt}
    &&\forall k\in K,\; r\in R,\; t\in T
    \label{fix:empty-link}
\end{align}
The constants are the tightest valid ones: no slot can hold more than
$\bar u^{V}_{n}=\lfloor W^{\max}/w_{n}\rfloor$ full RTIs of $n$ without
breaching the crush cap \eqref{sb:eq:stack-w-plant}, and no truck more
than $\bar u^{\mathrm{e}}_{r}=\lfloor pK^{w}/w^{\mathrm{f}}_{r}\rfloor$
folded RTIs of $r$ without breaching the payload
\eqref{sb:eq:truck-w-sup}, so neither row cuts a feasible load and
both are far smaller than the big-$M$ a generic gate would use.

In the integer feasible set of \emph{this} model the two rows are in
fact implied: every weight is strictly positive, so at $Y_{gkt}=0$ the
payload rows \eqref{sb:eq:truck-w-plant}--\eqref{sb:eq:truck-w-sup}
already force the truck's whole load to zero. They are nevertheless
written, for two reasons. \eqref{fix:slot-link} is strictly tighter than
that implication in the relaxation, because what bounds the content of
a single slot is the crush cap and not the truck payload. And the
implication is a property of the truck-indexed statement and of nothing
more general: the moment a layout is shared by several trucks --- the
restriction an implementation is tempted into, and the one solved in
Section \ref{sec:results} --- the capacity rows are written on the
shared layout while the payload row is written on the truck, so load
can be charged to a layout that no truck selected and
\eqref{sb:eq:agg-V}--\eqref{sb:eq:agg-Ve} count it as delivered. What
rules that out is gating the load on the truck's \emph{selection of the
layout it carries}. Here the two coincide, because a truck's layout is
its own, which is why \eqref{fix:slot-link}--\eqref{fix:empty-link} may
be written on $Y_{gkt}$; in the shared-catalogue restriction they must
be written on the layout-selection variable instead, and a gate on the
dispatch alone would leave the leak open.

The payload is a property of the truck:
\begin{align}
\sum_{n\in N} w_{n}\,\overline{V}_{knt}&\;\le\;pK^{w}\,Y_{pkt}
    &&\forall k\in K,\; t\in T
    \label{sb:eq:truck-w-plant}\\
\sum_{r\in R} w^{\mathrm{f}}_{r}\,V^{\mathrm{e}}_{krt}
    &\;\le\;pK^{w}\,Y_{skt}
    &&\forall k\in K,\; t\in T
    \label{sb:eq:truck-w-sup}
\end{align}
and the crush limit a property of each stack. On plant-bound trucks it
is period-exact, evaluated on the product mix actually loaded in the
slot; on supplier-bound trucks it is deterministic in the layout, since
the weight of a folded RTI depends only on its type:
\begin{align}
\sum_{n\in N} w_{n}\,V_{kint}&\;\le\;W^{\max}
    &&\forall k\in K,\; i\in I,\; t\in T
    \label{sb:eq:stack-w-plant}\\
\sum_{h\in H}\sum_{r\in R(h)}
    w^{\mathrm{f}}_{r}\,n_{hr}\,L_{s,i,h,k,t,r}&\;\le\;W^{\max}
    &&\forall i\in I,\; k\in K,\; t\in T
    \label{sb:eq:stack-w-sup}
\end{align}

% --------------------------
% Symmetry breaking and slot-set sizing
% --------------------------
\subsection{Symmetry breaking, slot-set sizing and data assumptions}
\label{sb:sec:symm}

The slots $i\in I$ of a given truck are interchangeable, so any
solution admits $|I|!$ permutations of the slot index. We break this
symmetry by enforcing that slot $i+1$ is used only if slot $i$ is used:
\begin{equation}
\widehat{X}_{g,i+1,k,t}\;\le\;\widehat{X}_{gikt}
\qquad \forall g\in G,\; k\in K,\; t\in T,\; i\in\{1,\dots,|I|-1\}.
\label{sb:eq:sym}
\end{equation}

The trucks of a given direction and period are interchangeable in the
same way: nothing in the data distinguishes one element of $K$ from
another. We likewise force them to be used in index order,
\begin{equation}
Y_{g,k+1,t}\;\le\;Y_{gkt}
\qquad \forall g\in G,\; t\in T,\; k\in\{1,\dots,|K|-1\}.
\label{dec:sym-truck}
\end{equation}
The row is valid because any solution can be re-indexed, but it removes
only the choice of \emph{which} subset of trucks is used. The loads and
the layouts remain interchangeable among the used trucks, so a plan
dispatching $m$ trucks in a period still appears $m!$ times in the
search, which is the difficulty Section \ref{sec:math} removes by
aggregation rather than by a row.

Five properties of the data are assumed throughout. The cardinality of
$I$ is a modelling parameter, and an upper bound that eliminates
provably unused slot variables without cutting any feasible layout is
\begin{equation}
|I|\;\le\;\left\lfloor
    \frac{pK\cdot qK}{\min_{h\in H} p_{h}\cdot q_{h}}
\right\rfloor,
\label{sb:eq:Ibound}
\end{equation}
corresponding to tiling the truck floor with the smallest-footprint
stack type; it serves as a tight preprocessing limit on $|I|$. The
product-to-RTI map is a partition --- each product travels in exactly
one RTI type, so $\rho$ is well defined and the families $M(r)$ are
disjoint --- and every RTI type carries at least one product, so no
$M(r)$ is empty. The lead time satisfies $lt\ge1$: a shipment never
arrives in the period it leaves, which is what the lagged terms of
\eqref{sb:eq:inv-plant-main} and \eqref{sb:eq:inve-sup-main} require.
All weights and heights are strictly positive, $w_{n}>0$,
$w^{\mathrm{f}}_{r}>0$ and $r_{gr}>0$; this is what makes
$\bar u^{V}_{n}$ and $\bar u^{\mathrm{e}}_{r}$ well defined, what lets
\eqref{sb:eq:height} force $L_{gihktr}=0$ at an empty slot, and what
makes \eqref{sb:eq:truck-w-plant}--\eqref{sb:eq:truck-w-sup} empty a
truck that is not dispatched.

Finally, the footprint geometry is consistent: one layer of stack type
$h$ tiles its footprint with RTIs of the type it carries,
$n_{hr}a_{r}=p_{h}q_{h}$ for every $h\in H$ and $r\in R(h)$, and the
volume charged to an RTI is its bounding box, $\gamma_{gr}=a_{r}r_{gr}$.
\emph{This is a departure}: the base formulation leaves $p_{h}$,
$q_{h}$, $n_{hr}$ and $\gamma_{gr}$ as four independent data items, and
without a relation between them the volume a layer occupies and the
volume its stack type reserves on the floor are unrelated numbers. Two
later arguments need the relation --- the inclusion of the 3D column
set in the volumetric one (Proposition \ref{prop:inclusion}) and the
floor-area inequality \eqref{bp:PP-areacut} that the pricing oracles
carry --- and both are false without it. It costs nothing here: the
instances of Section \ref{sec:instances} and the industrial catalogue
of Section \ref{sec:case} both build one stack type per distinct RTI
footprint, so $n_{hr}=1$ and the relation is an identity.

% --------------------------
% Domains
% --------------------------
\subsection*{Domains}

\begin{align*}
&Z_{gkt},\,Y_{gkt},\,\theta_{gkt},\,X_{gihkt},\,LX_{gihkt},\,LY_{gihkt},\,
  A_{gijkt},\,B_{gijkt},\,C_{gijkt},\,D_{gijkt}\in\{0,1\},\\
&L_{gihktr},\,V_{kint},\,V^{\mathrm{e}}_{krt},\,I_{gnt},\,Ie_{gtr}
  \in\mathbb{Z}_{\ge 0},\\
&PX_{gikt},\,PY_{gikt}\ge 0.
\end{align*}

% --------------------------
% The three modelling decisions
% --------------------------
\paragraph{Three decisions, and the alternatives they reject.}
Three of the rows above admit a tempting alternative that is wrong, and
each is argued where it occurs rather than left to the reader. A
fourth departure, on the data rather than on a row, is the footprint
consistency assumed in Section \ref{sb:sec:symm}. (i) The
utilisation floor is imposed on the load carried
\eqref{sb:eq:util-p}--\eqref{sb:eq:utilw-s} and not on the layer
capacity designed \eqref{sb:eq:util-cap}--\eqref{sb:eq:utilw-cap}; the
capacity form is a consequence and is kept as a tightening, never as
the requirement. (ii) The family weight $\omega_{pr}$ in that
tightening is the family \emph{maximum}: the minimum gives a strictly
stronger row that is not a necessary condition of a full truck, cuts
off feasible loads, and reports an optimum above the true one. (iii)
The configuration binary $Z_{gkt}$ and the dispatch binary $Y_{gkt}$
are both kept here, because the two differ --- at $Z=1,Y=0$ a truck
carries a layout and no freight --- but the difference is one the
reformulation of Section \ref{sec:math} does not need, and Section
\ref{sec:pp} deletes $Y$ from the pricing problem on an argument stated
there.

